\documentclass[12pt]{article}
\usepackage{amsmath,amssymb}
\usepackage{geometry}
\geometry{margin=1in}
\usepackage[hidelinks]{hyperref}
\setlength{\parindent}{0pt}
\usepackage{xcolor}
\usepackage{etoolbox}  %Supports toggle commands
\definecolor{darkblue}{RGB}{0,0,139}

%SETTINGS
\newtoggle{solutions} \togglefalse{solutions}

\begin{document}

\begin{center}
\textbf{Midterm Exam 1} \\
Introduction to Health Economics: EC 387 \\
Boston University \\
Spring 2026 \\
Professor Pauline Mourot
\end{center}

\vspace{1em}

\begin{itemize}
  \item The exam is scheduled for 75 minutes. There are a total of 100 possible points on the exam.  

  \item This exam contains 2 parts. Part I contains 7 true/false/uncertain questions. Part II contains 3 exercises.

  \item This is a closed-book, closed-note exam. You may use a calculator.

  \item If you use separate pieces of paper, please carefully number your answers so that it is clear which question you are answering.  

  \item \textbf{Please read the instructions carefully.}

  \item \textbf{You may receive partial credit if you show your work.} An incorrect answer with no justification or explanation will likely receive no credit.
\end{itemize}

\vspace{5mm}


\noindent \textbf{PART I: TRUE/FALSE/UNCERTAIN [35 points]}\\

For each of the questions below, please explain your answer.
The quality of the explanation is an important part of the grade for these problems and
can also help you earn substantial partial credit.
In answering each question, 1-3 sentences of explanation should be sufficient.

\vspace{7mm}

\noindent \textbf{QUESTION 1}\\
Research from a natural experiment in Oregon's Medicaid found that having health insurance decreases ER usage and increases the use of primary care physicians.

\iftoggle{solutions}{}{\vspace{50mm}
\newpage}

\iftoggle{solutions}{
{\color{darkblue}

\vspace{5mm}

FALSE / UNCERTAIN – The Oregon Health Experiment (Finkelstein et al. 2012, QJE) does find an increase in the use of primary care
but found no significant effect on inpatient admissions coming from the ER.
}}{}

\vspace{7mm}

\noindent \textbf{QUESTION 2}\\
Research from the RAND health insurance experiment and more recent research papers suggest that increases in cost-sharing reduce healthcare spending
through reducing mostly low-value care.

\iftoggle{solutions}{}{\vspace{50mm}}

\iftoggle{solutions}{
{\color{darkblue}

\vspace{5mm}

FALSE - The RAND health insurance experiment found that increases in cost-sharing reduced healthcare spending, but it reduced both highly effective and rarely effective care.
Brot-Goldberg et al (2024) also finds that both high-and low-value care are reduced when deductibles are higher.
}}{}

\vspace{7mm}

\noindent \textbf{QUESTION 3}\\
Employment in the healthcare industry has followed a similar trend as employment in manufacturing in the U.S. since 1990.

\iftoggle{solutions}{}{\vspace{50mm}}

\iftoggle{solutions}{
{\color{darkblue}

\vspace{5mm}

FALSE – If anything, it is the opposite. Employment in manufacturing has decreased since 1990 while it has steadily increased for healthcare.
See week 1 slide 2 for class 2.
}}{}

\vspace{7mm}


\noindent \textbf{QUESTION 4}\\
A risk-averse individual has a lottery ticket that would give them \$0 with 50\% probability and \$1,000 with 50\% probability.
They would prefer to take \$500 (for sure) instead of keeping the lottery ticket. 

\iftoggle{solutions}{}{\vspace{50mm}
\newpage}

\iftoggle{solutions}{
{\color{darkblue}

\vspace{5mm}

TRUE - The expected value of this lottery ticket is $\$0.5 \times 0 + \$0.5 \times 1000 = \$500$.
Getting \$500 for sure means receiving the utility from the expected value of the lottery ticket for sure.
With a concave utility function,
the utility from the expected value is above the expected utility from the lottery ticket,
so they will prefer to take \$500 for sure instead of keeping the lottery ticket.

}}{}

\vspace{7mm}

\noindent \textbf{QUESTION 5}\\
A risk-averse individual has a lottery ticket that would give them \$0 with 50\% probability and \$5,000 with 50\% probability.
They would prefer to take \$2,400 (for sure) instead of keeping the lottery ticket.

\iftoggle{solutions}{}{\vspace{50mm}}

\iftoggle{solutions}{
{\color{darkblue}

\vspace{5mm}

UNCERTAIN - Not enough information is provided.
The expected value of this lottery ticket is $\$0.5 \times 0 + \$0.5 \times 5000 = \$2500$.
The certainty equivalent of this lottery ticket is less than \$2500,
but we cannot determine whether it is above or below \$2400 without more information about the individual's utility function (degree of risk aversion).

}}{}

\vspace{7mm}

\noindent \textbf{QUESTION 6}\\
In the baseline Grossman model (1972) seen in class, the optimal health capital increases with age.

\iftoggle{solutions}{}{\vspace{50mm}}

\iftoggle{solutions}{
{\color{darkblue}

\vspace{5mm}

FALSE – An increase in age means higher age depreciation so a higher $\gamma$.
This means the user-cost of health is higher (rate of return).
We move along and up the MEC curve, corresponding to a lower optimal health capital.
See slide 21 of week 2 class 3.
}}{}

\vspace{7mm}

\noindent \textbf{QUESTION 7}\\
A health insurance plan with a premium of \$1,000 per year, a coinsurance rate of 35\%, no deductible and no out-of-pocket maximum
has an actuarial value of 35\%.

\iftoggle{solutions}{}{\vspace{40mm}
\newpage}

\iftoggle{solutions}{
{\color{darkblue}

\vspace{5mm}

FALSE – The actuarial value of this plan is $1 - 0.35 = 0.65$ or 65\%.
\begin{align*}
  AV &= \frac{\mathbb{E}[\text{Insurer Payments}]}{\mathbb{E}[\text{Total Spending}]}\\
  &= \frac{ \mathbb{E}[(1 - 0.35) * \text{Total Spending}]}{\mathbb{E}[\text{Total Spending}]}\\
  &= 0.65 \frac{\mathbb{E}[\text{Total Spending}]}{\mathbb{E}[\text{Total Spending}]}\\
  &= 0.65
\end{align*}
}}{}

\vspace{10mm}

% (See next page)
% \newpage

\noindent \textbf{PART II: Quantitative Exercises}\\

\vspace{5mm}

\noindent \textbf{Exercise 1 [25 points]}\\

Tony faces uncertainty about his health state next year.
There is a 80\% chance he will be perfectly healthy and have \$0 in medical bills,
but there is a 20\% chance he will get sick and have \$4,000 in medical bills.
His current wealth is \$12,000. His preferences are represented by the utility function

\begin{equation*}
U(w) = \sqrt{w}.
\end{equation*}

\medskip

\noindent \textbf{(a)} If he does not have insurance, what is Tony's expected wealth next year?

\iftoggle{solutions}{}{\vspace{40mm}}

\iftoggle{solutions}{
{\color{darkblue}
\begin{align*}
E[w] &= 0.8 \times 12,000 + 0.2 \times (12,000 - 4,000)\\
& = 11,200
\end{align*}
}}{}

\medskip

\noindent \textbf{(b)} If he does not have insurance, what is Tony's expected utility next year?
(round to the nearest decimal place)

\iftoggle{solutions}{}{\vspace{40mm}}

\iftoggle{solutions}{
{\color{darkblue}
\begin{align*}
E[U(w)] & = 0.8 \times \sqrt{12,000} + 0.2 \times \sqrt{12,000 - 4,000}\\
& = 105.5
\end{align*}
}}{}

\medskip

\noindent \textbf{(c)} If he does not have insurance,
what is Tony's certainty equivalent next year?
(round to the nearest decimal place)

\iftoggle{solutions}{}{\vspace{40mm}
\newpage}

\iftoggle{solutions}{
{\color{darkblue}
\begin{align*}
U(CE) &= E[U(w)] \\
\sqrt{CE} &= 105.5 \\
CE &= 11,130.3
\end{align*}
}}{}

\medskip

\textbf{(d)} Would Tony be willing to purchase a full insurance plan with a premium of \$2,000?

\iftoggle{solutions}{}{\vspace{40mm}}

\iftoggle{solutions}{
{\color{darkblue}
Expected utility if purchase is

\begin{equation*}
U = \sqrt{12,000 - 2,000} = 100 < 105.5,
\end{equation*}

so no.
}}{}

\medskip

\textbf{(e)} What would the actuarially fair premium for a full insurance plan for Tony be?

\iftoggle{solutions}{}{\vspace{40mm}}

\iftoggle{solutions}{
{\color{darkblue}

Actuarially fair price means $\text{Premium} = \mathbb{E}[\text{Cost}]$ (Profits of insurer are zero).
\begin{equation*}
Premium = 0.8 \times 0 + 0.2 \times 4000 = 800
\end{equation*}
}}{}

\vspace{5mm}

\noindent \textbf{Exercise 2 [20 points]}\\
Assume that an individual has income of \$80,000. The individual faces a risk of hospitalization,
which occurs with probability $p = 0.1$.
If this happens, the total cost will be \$15,000.
The individual's utility over income net of out-of-pocket medical costs is

\begin{equation*}
u(y) = \sqrt{y}.
\end{equation*}

Before facing the risk, the individual is deciding between two health insurance plans (Plan A and Plan B)
and remaining uninsured.

\medskip

\textbf{Plan A:}  
The individual pays a premium of \$1,050 (whether she shows up in the hospital or not) for a health insurance plan
with a deductible of \$5,000.
After paying out-of-pocket costs up to that deductible, the individual is fully insured.

\medskip

\textbf{Plan B:}  
The individual pays a premium of \$1,560 (whether she shows up in the hospital or not) for a health insurance plan
with no deductible and full insurance (so there are no out-of-pocket costs at all).

\medskip

\textbf{(a)} From the perspective of the firm, compute the expected cost and the expected profit to the insurer
from selling each plan.

\iftoggle{solutions}{}{\vspace{50mm}}

\iftoggle{solutions}{
{\color{darkblue}
Firm expected costs are how much they should expect to pay out. In the fully insured plan, this aligns with the consumer’s expected cost. In the plan with a deductible, the firm pays less than the full amount because of the deductible. Expected profits are the premiums minus the expected costs.

Plan A:

\begin{align*}
\text{Expected costs} &= 0.1 \times (15000 - 5000) = 1000 \\
\text{Expected profits} &= 1050 - 1000 = 50
\end{align*}

Plan B:

\begin{align*}
\text{Expected costs} &= 0.1 \times 15000 = 1500 \\
\text{Expected profits} &= 1560 - 1500 = 60
\end{align*}
}}{}

\medskip

\textbf{(b)} Does the consumer prefer to purchase Plan A, to purchase Plan B, or to remain uninsured?

\iftoggle{solutions}{}{\vspace{50mm}}

\iftoggle{solutions}{
{\color{darkblue}
To calculate consumer preferences, we use the utility function to calculate the consumer's expected utility
under each of the plans.

Uninsured:
\begin{equation*}
E[U] = 0.1 \times \sqrt{80000 - 15000} + 0.9 \times \sqrt{80000} = 280.0535
\end{equation*}

Plan A:
\begin{equation*}
E[U] = 0.1 \times \sqrt{80000 - 5000 - 1050} + 0.9 \times \sqrt{80000 - 1050} = 280.0761
\end{equation*}

Plan B:
\begin{equation*}
E[U] = \sqrt{80000 - 1560} = 280.0714
\end{equation*}

The consumer prefers to insure themselves. Further, they prefer Plan A to Plan B.
}}{}

\medskip

\textbf{(c)} Does the answer to part (b) change if the consumer is risk-neutral?

\iftoggle{solutions}{}{\vspace{40mm}}

\iftoggle{solutions}{
{\color{darkblue}
Now consider a risk-neutral utility function:

\begin{equation*}
u(y) = y.
\end{equation*}

Uninsured:
\begin{equation*}
E[U] = 0.1 \times (80000 - 15000) + 0.9 \times 80000 = 78500
\end{equation*}

Plan A:
\begin{equation*}
E[U] = 0.1 \times (80000 - 5000 - 1050) + 0.9 \times (80000 - 1050) = 78450
\end{equation*}

Plan B:
\begin{equation*}
E[U] = 80000 - 1560 = 78440
\end{equation*}

Now the consumer prefers being uninsured to paying for an insurance plan. Intuitive explanations here are also accepted.
}}{}

\medskip

\textbf{(d)} Suppose the government decides to make all high-deductible plans (such as Plan A) illegal.
If Plan A is eliminated, would the consumer prefer Plan B or prefer remaining uninsured?
Provide intuition for your answer in 1–3 sentences.

\iftoggle{solutions}{}{\vspace{40mm}}

\iftoggle{solutions}{
{\color{darkblue}
Using the same calculations from part (b), we can now see that consumers will prefer Plan B to being uninsured.
The consumer is still risk averse, so they are willing to pay for certainty rather than remaining uninsured.
Their level of risk aversion means that they are willing to still bear some risk and purchase a plan with a deductible,
but if that option is not available, they still prefer the less desirable insurance plan to being uninsured.
Their total willingness to pay is higher than the premium on Plan B.
}}{}


\vspace{10mm}



\noindent \textbf{Exercise 3 [20 points]}\\
This problem works through a simplified version of the Grossman model.
Individuals have the following utility function for consumption, $U(c)=c$.
Consumption is purchased from labor earnings, $y$, which is product of the individual's hourly wage 30 and total hours worked,
$T^w$. In other words,
\begin{equation*}
c = y = 30 \times T^w .
\end{equation*}

Total hours available is 24, and must be split between work hours $T^w$, hours spent sick $T^s$,
and hours spent producing health $T^H$.
Health, $H$, is produced solely from time spent producing health (e.g., exercise); i.e., $H = T^H$
and $T^S = 24 \times \exp(-H) = 24 \times \exp(-T^H)$.

\medskip

\noindent \textbf{(a)} What's the maximum utility the individual can achieve?\\
\textit{(Hint: $\log(24) \approx 3.18$.)}

\iftoggle{solutions}{}{\vspace{50mm}}

\iftoggle{solutions}{
{\color{darkblue}

\vspace{5mm}

\begin{align*}
c &= y = 30 \times T^w\\
& = 30 \times (24 - T^S - T^H) \\
  &= 30 \times (24 - 24 \exp(-T^H) - T^H)
\end{align*}

Taking the first-order condition with respect to $T^H$:

\begin{align*}
30 \times (24\exp(-T^H) - 1) &= 0 \\
\exp(-T^H) &= \frac{1}{24}\\
\frac{1}{\exp(T^H)} &= \frac{1}{24} \\
\exp(T^H) &= 24 \\
T^H &= \log(24)\\
T^H &\approx 3.18
\end{align*}

and 

\begin{align*}
    T^S &= 24 \exp(-T^H)\\
    &= 24 \frac{1}{\exp(\log(24))}\\
    & = 24 \times \frac{1}{24} \\
    &= 1
\end{align*}

and

\begin{align*}
T^w &= 24 - T^S - T^H \\
&= 19.82 \\
\end{align*}

and finally 
\begin{align*}
U(c) &= c = 30 \times T^w \\
&= 594.6
\end{align*}

}}{}

\medskip

\noindent \textbf{(b)} What's the allocation of time (i.e., $T^w$, $T^H$, and $T^S$) when the individual achieves the maximum utility level?

\iftoggle{solutions}{}{\vspace{45mm}}

\iftoggle{solutions}{
{\color{darkblue}

\vspace{5mm}

See derivations above:

\begin{align*}
T^S &= 1 \\
T^H &= \log(24) = 3.18 \\
T^w &= 24 - T^S - T^H = 19.82
\end{align*}

}}{}

\begin{center}
[END OF EXAM]
\end{center}

\end{document}